\subsection{TRANSPORTABLE RELATIONS}

Let \(\mathfrak{C}\) be a theory which is stronger than the theory of sets, let \(x_{1},\ldots ,x_{n},\) \(s_1,\ldots ,s_p\) be distinct letters which are distinct from the constants of \(\mathfrak{C}\) , and let \(\mathrm{A}_1,\ldots ,\mathrm{A}_m\) be terms in \(\mathfrak{C}\) in which none of the letters \(x_{i}\) \((1\leqslant i\leqslant n)\) and \(s_j\) \((1\leqslant j\leqslant p)\) appears. Let \(\mathrm{S}_1,\ldots ,\mathrm{S}_p\) be echelon construction schemes on \(n + m\) terms. Then the relation \(\mathrm{T}\{x_1,\dots,x_n,s_1,\dots,s_p\}\) :

\[
\begin{array}{c}
\text{``}s_{1} \in S_{1}(x_{1}, \ldots, x_{n}, A_{1}, \ldots, A_{m})
\text{ and }s_{2} \in S_{2}(x_{1}, \ldots, x_{n}, A_{1}, \ldots, A_{m}) \\
\text{and }\ldots\text{ and }s_{p} \in S_{p}(x_{1}, \ldots, x_{n}, A_{1}, \ldots, A_{m})\text{''}
\end{array}
\]

is called a typification of the letters \(s_{1}, \ldots, s_{p}\) .

Let \(\mathbb{R}\{x_1, \ldots, x_n, s_1, \ldots, s_p\}\) be a relation in \(\mathfrak{C}\) which contains certain of the letters \(x_i, s_j\) (and possibly other letters as well). Then \(\mathbb{R}\) is said to be transportable (in \(\mathfrak{C}\) ) with respect to the typification \(\mathbb{T}\) , the \(x_i (1 \leqslant i \leqslant n)\) being considered as principal base sets and the \(A_h (1 \leqslant h \leqslant m)\) as auxiliary base sets, if the following condition is satisfied: let \(y_1, \ldots, y_n, f_1, \ldots, f_n\) be distinct letters which are distinct from the \(x_i (1 \leqslant i \leqslant n)\) , the \(s_j (1 \leqslant j \leqslant p)\) , the constants of \(\mathfrak{C}\) , and all the letters which appear in \(\mathbb{R}\) or in the terms \(A_h (1 \leqslant h \leqslant m)\) , and let \(\operatorname{Id}_h (1 \leqslant h \leqslant m)\) denote the identity mapping of \(A_h\) on to itself. Then the relation

\begin{enumerate}
\def\labelenumi{(\arabic{enumi})}
\tightlist
\item
  ``T \(\{x_{1},\ldots,x_{n},s_{1},\ldots,s_{p}\}\) and ( \(f_{1}\) is a bijection of \(x_{1}\) onto \(y_{1}\) ) and \ldots{} and ( \(f_{n}\) is a bijection of \(x_{n}\) onto \(y_{n}\) )''
\end{enumerate}

implies, in \(\mathfrak{G}\) , the relation

\[
\mathrm{R} \left\{x _ {1}, \dots , x _ {n}, s _ {1}, \dots , s _ {p} \right\} \Longleftrightarrow \mathrm{R} \left\{y _ {1}, \dots , y _ {n}, s _ {1} ^ {\prime}, \dots , s _ {p} ^ {\prime} \right\},\tag{2}
\]

where

\[
s _ {j} ^ {\prime} = \left\langle f _ {1}, \dots , f _ {n}, \operatorname{Id} _ {1}, \dots , \operatorname{Id} _ {m} \right\rangle^ {S_j} (s _ {j}) \quad (1 \leqslant j \leqslant p).\tag{3}
\]

There is an analogous but simpler definition in the case where there is no auxiliary set.

For example, if \(n = p = 2\) and if the typification \(\mathbf{T}\) is ``\(s_1 \in x_1\) and \(s_2 \in x_1\)'', the relation \(s_1 = s_2\) is transportable. On the other hand, the relation \(x_1 = x_2\) is not transportable.

